\subsection{Distance Metrics}

Different distance metrics are used depending on the problem context. For two points $P = (P_x, P_y)$ and $Q = (Q_x, Q_y)$:

\subsubsection{Euclidean Distance}

The standard distance in the Euclidean plane:

\begin{equation}
    d_2(P, Q) = \sqrt{(P_x - Q_x)^2 + (P_y - Q_y)^2}
\end{equation}

This is the length of the straight line segment connecting the two points.

\subsubsection{Manhattan Distance}

Also known as taxicab distance or $L_1$ distance:

\begin{equation}
    d_1(P, Q) = |P_x - Q_x| + |P_y - Q_y|
\end{equation}

This represents the distance when movement is restricted to horizontal and vertical directions (like moving in a city grid).

\subsubsection{Chebyshev Distance}

Also known as chessboard distance or $L_\infty$ distance:

\begin{equation}
    d_\infty(P, Q) = \max(|P_x - Q_x|, |P_y - Q_y|)
\end{equation}

This represents the minimum number of moves a king needs to travel between two squares on a chessboard.

\subsubsection{Relationships Between Metrics}

For any two points $P$ and $Q$:

\begin{equation}
    d_\infty(P, Q) \leq d_2(P, Q) \leq d_1(P, Q) \leq 2 \cdot d_\infty(P, Q)
\end{equation}

\subsubsection{Coordinate Transformation}

Manhattan and Chebyshev distances can be converted through a 45° rotation:

\textbf{Manhattan $\rightarrow$ Chebyshev:}

Transform coordinates $(x, y) \mapsto (x + y, x - y)$

\textbf{Chebyshev $\rightarrow$ Manhattan:}

Transform coordinates $(x, y) \mapsto \left(\frac{x+y}{2}, \frac{x-y}{2}\right)$

This transformation is useful when a problem naturally uses one metric but is easier to solve with the other.
